% Part: lambda-calculus
% Chapter: syntax
% Section: terms

\documentclass[../../../include/open-logic-section]{subfiles}

\begin{document}

\olfileid{lam}{syn}{trm}
\olsection{الحدود}

مادة حدود حساب لامبدا متغيرات لا تنفد، هي
$\Obj{v_0}$ و$\Obj{v_1}$ و\dots، ورمز~«$\lambd$»، وأقواس؛ ومنها نبني الحدود استقرائيًا.
ونجعل $x$ و$y$ و$z$ و\dots{} أسماء للمتغيرات، و$M$ و$N$ و$P$ و\dots{} أسماء للحدود.

\begin{defn}[الحدود] \ollabel{defn:term}
هذه قواعد التعريف الاستقرائي لمجموعة \emph{حدود} حساب لامبدا:
\begin{enumerate}
  \item \ollabel{defn:term-var} كل متغير~$x$ يعدّ، بذاته، حدًا~$x$.
  \item \ollabel{defn:term-abs} إذا أخذنا متغيرًا~$x$ وحدًا~$M$، حصلنا على الحد
    $(\lambd[x][M])$.
  \item \ollabel{defn:term-app} من الحدين~$M$ و$N$ نحصل على الحد
    $(MN)$.
\end{enumerate}
\end{defn}

أما الحد~$(\lambd[x][M])$ الذي تبنيه القاعدة
\olref{defn:term-abs} فنسميه حاصل \emph{تجريد}؛ ونسمي~$x$ الوارد في
$\lambd[x]$ \emph{!!{parameter}}. وأما الحد~$(MN)$ الذي تبنيه
\olref{defn:term-app} فنسميه حاصل \emph{تطبيق}.

يلزم من هذا التعريف إثبات جميع الأقواس في الحدود؛ وقد تثقل بذلك الكتابة، كما ترى في الحد
$(\lambd[x][((\lambd[x][x])(\lambd[x][(xx)]))])$. وسنضع اصطلاحات تجيز إسقاط بعض الأقواس.
غير أن لهذا التعريف الرسمي فائدة، إذ يظهر منه بناء الحد بحسب
\olref{defn:term}. فالخطوة الأخيرة، مثلًا، في تكوين الحد
$(\lambd[x][((\lambd[x][x])(\lambd[x][(xx)]))])$ لا بد أن تكون تجريدًا فيه
!!{parameter} هو~$x$. فالحد ناشئ بالتجريد من
$((\lambd[x][x])(\lambd[x][(xx)]))$، وهذا ناشئ بتطبيق حدين، وكل منهما ناشئ بتجريد، ويستمر التحليل على هذا الوجه.

\begin{prob}
بيّن خطوات تكوين~$(\lambd[g][(\lambd[x][(g (x x))])
  (\lambd[x][(g (x x))])])$.
\end{prob}

\end{document}

